\documentclass[border=2pt]{standalone}
\usepackage{amsmath, amsthm, amsfonts}
\usepackage{tikz-cd, kotex}
\usepackage{quiver}
\usepackage{Operators}
\usepackage{palette}
\usepackage{diagram-style}
\usetikzlibrary{arrows.meta, positioning}
\begin{document}

%%FIG 두 almost split sequence 사이의 동형 (Auslander-Reiten_Theory-1.svg)
\begin{tikzcd}
	0 \arrow[r] & A \arrow[r, "f"] \arrow[d] & B \arrow[r, "g"] \arrow[d] & C \arrow[r] \arrow[d, equal] & 0 \\
	0 \arrow[r] & {A'} \arrow[r, "f'"'] & {B'} \arrow[r, "g'"'] & C \arrow[r] & 0
\end{tikzcd}

%%FIG A3 Auslander-Reiten quiver (Auslander-Reiten_Theory-2.svg)
\begin{tikzcd}[column sep=small, row sep=large]
	& & M_{[1,3]} \arrow[dr] & & \\
	& M_{[2,3]} \arrow[ur] \arrow[dr] & & M_{[1,2]} \arrow[dr] & \\
	M_{[3,3]} \arrow[ur] & & M_{[2,2]} \arrow[ur] & & M_{[1,1]}
\end{tikzcd}

\end{document}
